% 本文件由 LaTeX 排版 Agent 根据有效 translation.json 自动生成。
% author=Peter of Alexandria; work_id=0619; title=亚历山大城的彼得行传

\chapter{\WorkTitle{亚历山大城的彼得行传}}

% structure: p0001 source=h1 target=omitted reason=page-title-already-rendered-by-parent

若我全身的肢体都变成舌头，所有肢体的关节都发出清晰的声音，也远不足以表达我们至可敬的教父、亚历山大城总主教彼得是谁、何其伟大、何其良善。尤其我认为，将他在暴君手下所受的艰危、他与外邦人和异端者所经历的争战诉诸笔墨，实属不妥，免得我似乎将这些作为我颂词的主题，而非他那为了保全天主之民而勇敢承受的苦难。然而，由于叙述者之职在叙述他内心的交往和奇事时必定力有未逮，语言也绝不足以为此任务，我以为只宜描述他那些已知使他获得主教高位的事迹，以及将亚略从教会合一中剪除后，他以殉道者的桂冠受冕的经过。然而，我认为这是一个光荣的结局，是一场伟大竞赛的景象，足以令那些不怀疑真实叙述、且叙述未被谎言玷污的人满足。因此，在开始叙述这位至圣之人的主教任期时，让我们求助于他自己的言语，以便使之与我们的文风相配合。\par

亚历山大城是一座极其宏伟的城邑，不仅在埃及人中居首位，在底比斯人和距埃及不远的利比亚人中也是如此。从我们的主和救主耶稣基督降生成人以来，二百八十五年一个周期已经过去，当时可敬的德奥纳，本城的主教，以超凡的飞翔升入天国。彼得继他之后，执掌教会的船舵，由全体圣职人员和整个基督徒团体任命为主教，是从宗徒圣史玛尔谷（他也是本城的总主教）以来的第十六位。他确实如同晨星升起在群星之中，以神圣德行的光辉照耀，极其宏伟地治理了信德的堡垒。他在圣经知识上不逊于任何前任，卓越地为教会的利益和训诲而努力；他具有非凡的明智，事事完美，是天主的真司铎和祭品，日夜警醒劳苦于每一项司祭的职责。\par

但因德性是热忱者的标志，\Quote{高山之巅易遭雷击}，由此他遭受了与对手的多种冲突。何须多言？他几乎一生都在迫害中度过。同时，他祝圣了五十五位主教。最后，墨利提乌斯——在心志和名号上最为黑暗——成了吕科波利斯城分裂派的主教，做了许多违反教规之事，甚至在残暴上超过了那些在主受难时不敢撕裂祂长衣的嗜血兵士；他因放纵疯狂而被驱策，以至于撕裂了天主教会——不仅在埃及诸城，甚至在村庄里——按自己的派别祝圣主教，毫不关心彼得，也不关心在彼得身上的基督。亚略当时还是个平信徒，尚未受圣职剃发，依附于他，并对他及其家族极其亲爱；这并非无因：正如圣经所说，\Quote{每种动物都喜爱自己的同类}。但这事被天主的人得知后，他深感忧伤，说这迫害比先前的更甚。虽然他正在躲避，但仍尽其所能，到处发出劝勉，宣讲教会的合一，坚固人们对抗墨利提乌斯的无知和邪恶的鲁莽。因此，有不少人受他有益劝诫的影响，离开了墨利提乌斯的邪恶。\par

几乎与此同时，亚略武装以毒蛇的狡诈，仿佛离开墨利提乌斯一方，逃到彼得那里避难。彼得应主教们的请求，提升他至执事尊位，却不知他极端的伪善。因为他就像一条充满致命毒液的蛇。然而，对这虚伪者行覆手礼也不能归罪于这位圣人，正如西满的幻术不能归咎于斐理伯一样。同时，墨利提乌斯的可憎邪恶大大增长；有福的彼得担心异端瘟疫蔓延至托付给他照管的整个羊群，又知光明与黑暗没有相通，基督与贝里雅耳没有和谐，便以书信将墨利提乌斯派从教会共融中分离。因邪恶的性情不能长久隐藏，邪恶的亚略见他的支持者和同谋被从教会尊严中推翻，即时陷入悲伤和哀叹。这并未逃过圣人的眼睛。因为当他的伪善被揭露时，立即使用福音之剑：\Quote{若你的右眼使你跌倒，就挖出它，从你身上扔掉}，\BibleRef{玛 5:29}，并将亚略如同腐坏的肢体从教会身上割除，驱逐并逐出信友的共融。\par

此事之后，迫害的风暴突然减弱，和平虽短暂却微笑降临。那时，主这位最精选的司铎在人民面前明显闪耀，信友们开始成群奔去纪念殉道者，并聚集在会众中赞美基督。这位神圣法律的司铎以神圣的言词激励他们，如此激发和坚固他们，以致信徒的数目在教会中不断增长。但人类救恩的古老仇敌并未长久安静并以恩待的眼光看待这些事。因为忽然间，异教的暴风发出敌对的雷声，如同冬日的骤雨冲击教会的宁静，并将其驱散逃亡。但为更清楚地理解此事，我们必须回溯到那不敬神、背叛天主的迪奥克勒提安，以及同时与其子马克西敏用暴政骚扰东方地区的马克西米安·加莱里乌斯的暴行。\par

因为在此人之时，基督徒的迫害之火如此猛烈，以至于不仅在宇宙的一隅，甚至在整个世界，无论陆地和海洋，不敬之风雷声大作。御旨和最残酷的法令四处传布，基督的敬拜者时而公开被杀，时而遭到暗算；没有一天一夜不流基督徒的血。杀戮的方式也非止一种：有些人被各种最痛苦的酷刑处死；有些人则被剥夺亲人的关怀和故土的安葬，被流放到异地，通过某种前所未闻的新刑罚手段被驱向殉道的终点。哦，何其可怕的邪恶！他们的不敬如此之大，甚至将神圣敬拜的圣所连根拔起，并将圣书投入火中。那可咒诅的记忆中的\Person{戴克里先}死后，\Person{君士坦丁大帝}被选为治理王国的君主，并在西方地区开始执掌政权。\par

那些日子，有人将上述总主教的消息禀报给\Person{马克西米努斯}，说他是一名领袖且在基督徒中居首位；此人因惯常的邪恶而激怒，当场下令逮捕\Person{伯多禄}并投入监狱。为此，他派遣了五位\OriginalTerm{护民官}{tribunes}，随同其士兵队伍前往\Place{亚历山大}；他们奉命来到后，突然抓住了基督的司铎，将他囚禁在狱中。信众的虔敬何等奇妙！当得知这位圣人被关在监狱的地牢中时，一大群人——主要是隐修士和贞女团体——蜂拥而至，他们不带任何武装，却以泪河和虔诚的慈爱围住了监狱的四周。如同孝顺的儿子对待慈父，更确切地说，如同基督徒的肢体对待基督的元首，他们以满腔的同情紧紧依附着圣人，成为他的围墙，使任何外邦人无法接近他。众人的心愿、呼声、同情和决心完全一致：宁愿死也不愿看到这位圣人遭受任何不幸。天主的人被囚禁在同样的脚镣中数日，身体被向后推压；护民官就此向国王提出建议，国王按其残暴的作风，判处这位最蒙福的宗主教死刑。此事传到基督徒耳中，他们同心合意开始以哀叹和哭号守护监狱的入口，坚决阻止任何外邦人接近他。当护民官无论如何都无法接近他以执行死刑时，他们商议决定：士兵们拔剑冲入人群，将他拉出来斩首；若有人阻挡，就当场处死。\par

与此同时，\Person{亚略}尚只被授予了肋未人的品位，他害怕在如此伟大的教父去世后，自己再也无法与教会和好，于是来到那些在神职人员中居首的人那里；他假意悲叹，以哀恳的话语和动听的言辞，试图说服圣总主教向他施以怜悯，解除他的绝罚。然而，有什么比伪装的心更欺骗人呢？有什么比圣洁的镇静更单纯呢？毫不延误：被请求的人进入基督的司铎那里，照常致词后，俯伏在地，以哀叹和眼泪亲吻他的圣手，恳求他说：\Quote{最蒙福的教父，因你信德之卓越，主已召唤你去接受殉道的荣冠，我们毫不怀疑这荣冠很快就要临于你。因此，我们认为你应本着你惯常的虔敬，宽恕\Person{亚略}，垂顾他的哀吁而施予宽容。}\par

天主的人听了这话，义愤填膺，推开他们，举手向天，喊道：\Quote{你们竟敢为\Person{亚略}向我求情？\Person{亚略}，无论在此世还是来世，将永远被驱逐，与天主子、我主耶稣基督的荣耀隔绝。}他这样抗议，在场众人都惊惧万分，如哑巴一般，默然不语。他们更怀疑他并非没有某种神的启示，才发出对\Person{亚略}的判决。但当仁慈的父亲看到他们因良心自责而沉默忧伤，便不愿固执于严厉，也不愿好像轻视他们而不予安慰；于是他叫来在司铎中看来较年长、最圣洁的\Person{阿基拉斯}和\Person{亚历山大}，让他们一左一右，稍离众人，在讲话结束时对他们说：\Quote{弟兄们，不要以为我是个不人道、严厉的人；因为我自己也生活在罪的律法之下。但要相信我的话：\Person{亚略}隐藏的诡计超越一切不义和邪恶，我并非自作主张批准了他的绝罚。因为今夜，当我庄严地向天主倾心祈祷时，一个约十二岁的男孩站在我面前，他的面容光辉得使我无法承受，因为我们所在的整个小屋被大光照耀。他穿着亚麻长袍，从颈到脚分为两片，双手拿着长袍的裂缝，贴在胸前以遮盖裸体。见此异象，我惊愕不已。当我被赐予说话的勇气时，我喊道：主，谁撕裂了你的长袍？他回答说：\Person{亚略}撕裂了它，你千万要小心，不可接纳他共融；看哪，明天他们会来为你求情。所以，务必不要被说服同意；反之，你要命令\Person{阿基拉斯}和\Person{亚历山大}司铎，他们将在你离世后管理我的教会，绝不可接纳他。你很快就会完成殉道者的份。这异象没有别的原因。现在我已让你们满意，我已向你们宣告了我所受的命令。但你们将因此如何行，是你们自己的事。}\\newline{}关于\Person{亚略}，就说到这里。\par

\Emph{他继续说：}亲爱的，你们也知道，且深知我在你们中间的生活是怎样的，以及我从那些不认识主和救主、疯狂宣扬众多并非神的偶像的外邦人中忍受了多少争战。你们也知道，为躲避迫害者的怒气，我如何漂泊流徙，从一处到另一处。我长期隐藏在\Place{美索不达米亚}，也在腓尼基人中的\Place{叙利亚}；在\Place{巴勒斯坦}各地也流浪了很久；从那里，若可以这样说，在另一种气候中，即岛屿上，我停留了不短的时间。然而，在这些灾难中，我日夜不停地写信给我所受托的可怜的羊群，在基督的合一中坚固他们。因为对他们焦虑的挂虑不断催促我的心，不让我安息；只有当我将他们交托给至高之权时，我才觉得稍可忍受。\par

同样，因那些有福的主教们——我指\Person{斐禄斯}、\Person{赫西基乌斯}和\Person{狄奥多鲁斯}——领受了相称的神圣恩宠的召叫，我的心灵经历了何等大的激动！因为这些人，如你们所知，为基督的信德与其他宣信者一起，忍受了各种酷刑。又因为在这场斗争中，他们不仅为圣职界，也为平信徒作了旗手和导师，我为此极其害怕，唯恐他们在长期磨难中跌倒，也怕他们的背弃（说起来可怕）成为许多人绊倒和否认信德的起因，因为与她们一同被囚在监狱围墙内的有六百六十多人。因此，虽然劳苦沉重，我不断写信给他们，提及所有那些预言过的经文，劝勉他们以神圣灵感的力量赢得殉道者的棕榈枝。但当我听说他们卓越的坚忍以及他们所有人光荣的苦难结局时，我俯伏在地钦崇基督的尊严，祂认为配得将他们列在殉道者的行列中。\par

我何必对你们提起\Place{吕科波利斯}的\Person{梅肋提乌斯}？他对我施加了何等迫害、何等诡计，我相信你们都很清楚。哦，可怕的邪恶！他竟敢撕裂圣教会，这教会是天主子用祂的宝血赎回的，并且为了将她从魔鬼的暴政中解救出来，祂毫不迟疑地舍弃了生命。这教会，正如我开始说的，邪恶的\Person{梅肋提乌斯}撕裂它，不停地将其监禁在牢狱中，甚至折磨圣主教们——他们比我们先一步经由殉道进入了天堂。因此，要警惕他的阴谋。因为我，如你们所见，被神圣的爱德束缚，视天主的旨意高于一切。我确实知道，护民官们低声急切地密谋我的死亡；但我不会因此与他们有任何沟通，也不会视我的生命比我自己更宝贵。相反，我准备完成我主耶稣基督曾恩许我要跑完的赛程，并忠实地向祂交还我从祂那里接受的职务。请为我祈祷，弟兄们；你们将不再看到我在此生与你们同住。因此，我在天主和你们的弟兄团体前作证，我在你们所有人面前保全了清洁的良心。因为我未曾避讳向你们宣布主的训令，也未曾拒绝向你们宣示将来必要的事。\par

因此，你们要谨慎，并顾及全群，圣神立你们为监督，继承我的职位——首先是\Person{Achillas}，其次是你\Person{Alexander}。看，我用活生生的声音向你们证明：在我死后，教会内将有人起来讲说邪说，又会像\Person{Meletius}那样再次分裂教会，随从他们之便勾引人众。所以，我事先已告诉你们。但我恳求你们，我自己的心肠，要警醒；因为你们必须经历许多磨难。我们并不比我们的祖先更好。你们不知道我的父亲——那养育我的至圣主教\Person{Theonas}——从外邦人那里遭受了什么吗？他的主教座我承继了。但愿我也拥有他的品行！为何还要提他的前任伟大的\Person{Dionysius}呢？他四处漂泊，从疯狂的\Person{Sabellius}那里遭受了许多灾祸。我也不愿忽略你们，至圣的父老和神圣法律的最高司祭\Person{Heraclius}和\Person{Demetrius}，\Person{Origen}——那邪说制造者——曾为他们设下许多诱惑，他在教会中引发了可恶的分裂，这分裂至今仍使教会混乱。但那时保护他们的天主的恩宠，我相信也会保护你们。可是，我极明达的弟兄们，为何我要用冗长的言词拖延你们呢？剩下的是，我以那位如此祈祷的宗徒最后的话向你们致辞：\Quote{现在我把你们交托给天主和祂恩宠的圣言，这圣言有能力引导你们和祂的羊群。}\BibleRef{宗 20:32}说完这话，他跪下同他们祈祷。话毕，\Person{Achillas}和\Person{Alexander}亲吻他的手和脚，放声痛哭，尤其为他所说的话悲伤，因为他曾说他们今后在今生不会再见到他。然后，这位最温和的老师走向其余的神职人员——如我所说，他们曾为\Person{Arius}的事进来与他交谈——向他们说了最后安慰的话和必要的话；接着向天主倾心祈祷，向他们告别，平和地遣散了众人。\par

事情如此结束后，各处纷纷传扬，\Person{Arius}被从公教会的合一中开除，实有天主之介入。但那诡计多端者，万恶的散布者，并未停止将蛇毒藏在他心腹的迷宫中，希望藉\Person{Achillas}和\Person{Alexander}得以和好。这就是异端魁首\Person{Arius}，分裂同质且不可分的天主圣三者。就是他，用狂妄邪恶的口，不怕凌辱主和救主，超过一切其他异端者——我说的是主和救主，祂因怜悯我们人类的迷途，又因悲悯世界将死于致命的毁灭与定罪中，为我们众人屈尊在肉身内受难。因为不可相信那不能受苦的天主性会受难。但由于神学家和教父们已以更好的方式将此类亵渎从公教徒的耳中除去，而我们的任务另有其所在，就让我们回到主题吧。\par

这位极其智慧的大司牧觉察到护民官的残酷计谋，他们为了置他于死地，不惜刀剑加于在场的全体基督徒群众。他不愿他们与他同尝死亡的苦味，而是如同忠信的仆人效法他的主和救主——祂的行为与言语一致：\BibleQuote{善牧为羊舍命}\BibleRef{若 10:11}——他出于虔敬，召叫一位在场听他训诲的长老，对他说：\Quote{你去对那些要杀我的护民官说：停止你们的一切焦虑，看！我自愿并乐意将自己交给他们。叫他们今夜到这监狱屋子的后面，在他们听见墙内有信号的地方，就挖掘，把我带走，随他们奉命而行。}那位长老服从了这位至圣之人的命令——因为不能违抗如此伟大的父老——就去了护民官那里，并按吩咐告诉了他们。他们听了后非常欢喜，带着一些石匠，大约黎明时分，不带士兵，来到指定地点。天主的人整夜不眠，持守祈祷与警醒。但当他听见他们走近时，所有与他在一起的人都沉睡了，他以缓慢轻柔的脚步下到监狱内部，照约定敲了墙壁；外面的人听见后，挖开一个洞口，接收了这位基督的勇士——他全身武装，没有铜甲，却有主十字架的德能，并准备好实践主的话：\BibleQuote{你们不要害怕那杀害肉身而不能杀害灵魂的；但更要害怕那能把灵魂和身体都投入地狱中的。}\BibleRef{玛 10:28}奇妙的事发生了！那夜风雨极大，以致看守监狱大门的人没有一人听见挖掘的声音。这位极其坚毅的殉道者还催促那些凶手说：\Quote{趁那些看守我的人尚未察觉，快做你们要做的事。}\par

但他们将他抬起，带到称为布科利亚的地方，圣马尔谷在那里为基督殉道。圣人们的德能何等奇妙！当他们抬着他前进，看到他在面临死亡危险时如此坚定和刚毅，突然恐惧和战栗降临到他们身上，以致无人能注视他的面容。此外，这位蒙福的殉道者请求他们允许他到圣马尔谷的墓地去，因为他想将自己托付于他的庇佑。但他们因慌乱而低头看着地面，说：\Quote{你愿怎样就怎样，但要快些。}于是他走近福音作者的安葬之处，拥抱它，并向他说话，仿佛他还活在肉身中并能听见一样，如此祈祷：\Quote{极可敬的父啊，你是独生救主的福音作者，你见证了他的苦难，基督——我们众人的救主——拣选了你作为此座的第一位主教和支柱；他交托给你向整个埃及及其边界宣讲信德的任务。我所说的你，忠实地完成了那交托给你的、为人类救恩服务的职分；作为这劳苦的酬报，你无疑已获得了殉道者的棕榈枝。因此，你不仅正义地被尊称为福音作者和主教。你的继任者是阿尼亚努斯，其余的后继者依次而下，直到最蒙福的德奥纳斯，他管教了我的幼年，并屈尊教育了我的心。我，一个罪人和不配的人，超过我应得的，通过继承的传承被指定为他的继任者。而最美好的是，看！神圣恩赐的丰盈已经赐予我成为祂宝贵十字架和喜乐复活的殉道者，给予我的虔敬以祂苦难的甘甜宜人的馨香，使我得以向祂倾流我的血作为祭献。因这祭献的时刻已近，请为我祈祷，好使神圣的德能助我，使我以刚强的心和坚定的信德，配得此苦战的终点。我也将托付给我牧养的信众团体，即基督的敬拜者，交托于你光荣的庇佑；我向你说，我以祈祷交托他们，你被认可为所有此前和此后占据此教座者的创始者和守护者，你持有此座最初的光荣，并非人的继任者，而是神-人基督耶稣的继任者。}说完这些话，他从圣墓稍后退，举起双手向天，大声祈祷说：\Quote{啊，独生的耶稣基督，永恒圣父的圣言，求祢俯听我呼求祢的仁慈。我恳求祢，向祢教会中的风暴说平安，并以我——祢仆人——的血的倾流，使祢子民的迫害终结。}这时，一位献身于天主的童贞女，她的独修小室毗邻福音作者的坟墓，当她彻夜祈祷时，听见从天上有声音说：\Quote{伯多禄是宗徒中的第一位，伯多禄是亚历山大里亚殉道主教中的最后一位。}\par

他结束祈祷后，亲吻了真福福音作者和那里埋葬的其他教长的坟墓，然后走向护民官。但他们看到他的面容如同天使的面容，惊恐畏惧，不敢向他提及他迫在眉睫的苦战。尽管如此，由于天主不遗弃那些信赖祂的人，祂不愿在如此大考验的时刻使祂的殉道者失去安慰。因为，看！一位老人和一位年迈的童贞女，从小镇赶来向城市前进，其中一人携带四张兽皮兜售，另一人带着两床亚麻布。蒙福的主教看见他们，辨认出这是天主针对他自己的神圣安排。他立刻问他们：\Quote{你们是基督徒吗？}他们回答说：\Quote{是。}于是他说道：\Quote{你们要去哪里？}他们回答说：\Quote{到城里的市场去卖我们携带的这些物品。}于是这位最仁慈的父回答说：\Quote{我忠信的孩子们，天主已拣选了你们，与我一同坚持吧。}他们立刻认出他来，说：\Quote{大人，愿照您所吩咐的而行。}然后他转向护民官，说：\Quote{来吧，做你们要做的，完成王的命令；因为天快亮了。}但他们，仿佛因王子邪恶的法令而被迫，将他带到福音作者圣所对面的一处地方，靠近坟墓的一个山谷。于是圣人说：\Quote{老人，铺开你携带的兽皮；老妇人，也铺开亚麻布。}他们铺开后，这位最坚定的殉道者登上去，伸出双手向天，双膝跪地，心智凝视天空，向这场竞赛的全能审判官献上感恩，并用十字圣号坚固自己，说：阿们。然后他从颈上解下\OriginalTerm{肩带}{omophorion}，伸出来，说：\Quote{命令你们的，快做吧。}\par

与此同时，护民官的手瘫痪了，他们彼此相看，每个人都催促同伴动手，但他们全都被惊愕和恐惧所抓住。最后他们商定，从公款中拨出一笔赏钱给行刑者，并任由敢于犯下谋杀的人享有这笔赏钱。他们毫不迟疑，每人拿出五枚\OriginalTerm{索利多}{solidi}。但正如外教诗人所说——\par

\Quote{你这可诅咒的黄金欲，人心还有什么不能驱使？} 其中一人，效法叛徒犹达斯，被金钱欲所怂恿，拔出剑来，斩杀了主教，于十一月二十五日，在他担任主教十二年之后——其中三年是在迫害之前，其余九年则是在各种迫害中度过的。刽子手立即索取血金，这些邪恶的买主，更确切地说是生命的毁灭者，迅速返回，因为他们害怕民众的众多，因为正如我所说，他们没有军事护卫。但真福殉道者的遗体，如最早前往刑场的教父们所证实的，仍然直立着，仿佛正在祈祷，直到许多人聚集前来，发现它保持同一姿势；因此，他生前的惯常行为，由这无生命的身体作证。他们还发现一位老年男子和老年妇女，满怀悲伤和哀叹，看守着教会最珍贵的遗物。于是，他们以凯旋式的葬礼尊崇他，用亚麻布遮盖他的身体；而他所倾流的圣血，他们虔诚地收集在一个袋子里。\par

与此同时，从人口稠密的城市蜂拥而至的无数男女，呻吟着、呼喊着，彼此询问事情是如何发生的，因为他们不知道真相。事实上，从最小的到最大的，所有人心中都弥漫着巨大的悲伤。因为当城中的首领们看到民众可嘉的坚持——他们正忙于分取他的圣物作为遗物——他们便用兽皮和亚麻布将他裹得更紧。因为至圣的天主仆役总是穿着白色的司祭服饰，即长袍、\OriginalTerm{科洛比翁（短袖衣）}{kolobion} 和 \OriginalTerm{奥莫福里翁（肩带）}{omophorion}。于是他们之间起了不小的争执；有些人想把至圣的肢体运到他亲自建造、如今安息的教堂去；另一些人则尽力将他运到圣史堂，即他获得殉道终点的地方。既然双方互不相让，他们便开始将虔诚的敬礼转变为争吵和权利之争。与此同时，一群从事公共运输服务的勇敢的元老们，看到所发生的事——因为他们就在海边——备好一艘船，突然抓起圣髑，放在船里，然后从后面攀上\Place{法罗斯}，经过一个名叫\Place{勒乌卡多}的区域，来到了至真福\Person{天主之母}、终身童贞\Person{玛利亚}的教堂，正如我们开始所说的，这座教堂是他建在西郊的郊区，作为殉道者的墓地。于是，人群仿佛天堂的宝藏被夺走了一般，有些人沿着直路，另一些人则绕行更曲折的路，匆忙跟随。当他们终于到达那里时，对于将他安放在何处已无任何争议，而是通过共同无可非议的商议，他们同意先将他放在他的主教座上，然后再埋葬他。\par

最谨慎的读者，我不会让你把这当作狂想和迷信，因为如果你了解这新奇之事的缘由，你会钦佩和赞许民众的热忱和行为。这位真福司铎，当他举行神圣奥迹的圣事时，并不像教会的习惯那样坐在他的主教宝座上，而是坐在其下的脚凳上。人民看到后，表示不满，并抱怨地喊道：\Quote{神父，您应当坐在您的椅子上。} 当他们反复这样说时，主的仆役站起来，以平静的声音平息了他们的抱怨，再次坐在同一脚凳上。所以这一切似乎都是出于谦逊的动机而做。但在某个大节日，他正在奉献弥撒祭献，想要做同样的事。于是，不仅人民，连圣职人员也异口同声地喊道：\Quote{主教，请坐在您的椅子上。} 但他仿佛意识到某个奥秘，拒绝这样做；他做出静默的手势——因为没有人敢固执地违抗他——他使他们全都安静下来，然而，他仍然坐在椅子的脚凳上。弥撒的庄严礼仪照常举行后，每位信友各自回家。\par

但是，天主的人召来圣职班，以平静安祥的心志，指责他们行事鲁莽，说：\Quote{你们怎么不因附和百姓的呼喊并责备我而羞愧呢？然而，既然你们的责备不是来自骄傲的浑浊激流，而是来自纯洁的爱泉，我就向你们揭示这奥秘的秘密。每当我想要靠近那座位时，我就看见一种德能仿佛坐在上面，极其光辉灿烂。于是我在喜乐与恐惧之间悬置，承认我完全不配坐在这样的座位上，若我不顾虑会给百姓造成冒犯的因由，无疑我甚至不敢坐到那椅子本身。所以，我可爱的儿子们，在这件事上，我似乎在你们面前违反了主教规则。尽管如此，许多次当我看它空着时，正如你们自己所见证的，我并不拒绝照常坐到座位上。因此，你们现在既已得知我的秘密，并且确信若我蒙准许，我会坐到座位上，因为我并非轻视我职位的尊严，就要停止再参与百姓的呼喊。}这位至圣之父在生前不得不向圣职班作出这一解释。因此，基督的信徒怀着虔敬的记忆这一切，将他的圣躯抬来，使其坐在主教宝座上。那时，民众中升起了如同欢送他仍活着在身体中的那般巨大的喜乐和狂喜。然后，他们用芳香的香料为他敷抹，用丝绸裹包；他们中的每一位，凡能最先带来的，都视为最大的收获。然后，他们手持象征胜利的棕榈枝，燃着明烛，高唱赞歌，焚香芬芳，庆祝他天上胜利的凯旋，安放了圣髑，并将他埋葬在他很早以前建造的墓园里；从那时起，直到今天，神奇的德能从未停止显现。诚然，虔诚的誓愿蒙受垂听赐福；弱者的健康得以恢复；不洁之神的驱逐见证着殉道者的功德。这些恩赐，主耶稣啊，都是祢的，祢惯于如此豪迈地荣耀祢的殉道者于死后：祢与圣父及\OriginalTerm{同体圣神}{Holy Consubstantial Spirit}，永生永王，万世无疆。阿们。此后，那豺狼和诡计的制造者，即\Person{亚略}，披着羊皮进入主的羊栈扰乱折磨羊群，或他如何得以升到司祭品位，让我们简要叙述。这并不是要烦扰那些敢于将已被天上的扇子从教会中簸出去的背道和传染的莠子带回主禾场上的人，因为这些人无疑被视为圣德卓越，但他们以为相信这样一位圣人是小事，于是违背了神命的诫命。那么怎样呢？我们谴责他们吗？绝不。因为只要这朽坏的身体压迫着我们，这地上的居所压抑着我们软弱的知觉，许多人就容易被自己的臆想所欺骗，认为不义的是正义的，不洁的是圣洁的。\Person{基贝红人}，虽按神的警告本应尽行毁灭，他们心里所愿与口中和神色不同，却能够迅速欺骗了\Person{若苏厄}——那位公义的应许之地的分配者。\Person{达味}也充满先知灵感，当他听到那欺骗青年的言语时，虽然是出于天主深奥且公义的审判，但行事却与事情的真实情况大相径庭。还有什么比宗徒们更高超呢？他们并没有脱离我们的软弱。因为其中一位写道：\Quote{我们众人都犯许多过失。}\BibleRef{雅 3:2} 另一位写道：\Quote{如果我们说我们没有罪，便是欺骗自己，真理不在我们内。}\BibleRef{若一 1:8} 但当我们为这些悔改时，就更容易获得宽恕，因为我们不是故意犯罪，而是出于无知或软弱。当然，这类过失不是出于偏离正道，而是由于慈悲的宽容。但我把为此辩护的事留给他人去做；让我们继续眼前的事。在那位卓越的信德护卫者、名副其实的\Person{伯多禄}，借着殉道的胜利之后，等等。\par

\EditorialNote{其余部分缺失。}\par

\SourceRecord{Translated by James B.H. Hawkins.；Ante-Nicene Fathers；Vol. 6.；Edited by Alexander Roberts, James Donaldson, and A. Cleveland Coxe.；Buffalo, NY: Christian Literature Publishing Co.,；1886.；<http://www.newadvent.org/fathers/0619.htm>.}
